% =====================================================================================
% Instrument validity. Included from kernelascent_full.tex before the Synthesis, because it
% qualifies every null the Results section reports.
% =====================================================================================
A benchmark that reports a null owes its reader a second result: evidence that the
measurement \emph{could have} reported a positive. Self-improvement benchmarks are unusually
exposed here, because the headline quantity is a difference between two arms of the same
pipeline. Any defect that suppresses both arms equally leaves the contrast at zero while every
surface diagnostic --- jobs complete, kernels compile, scores populate --- looks healthy. A
broken instrument does not report that it is broken; it reports $\Delta \approx 0$.

Over this project, \textbf{five independent pipeline defects each produced a plausible,
publishable-looking null}, summarised in Table~\ref{tab:defects}. None announced itself, and
three survived a code review before being caught by an assertion.

\begin{table}[t]
\centering\small
\begin{tabular}{@{}p{0.30\linewidth}p{0.31\linewidth}p{0.31\linewidth}@{}}
\toprule
\textbf{Defect} & \textbf{Presented as} & \textbf{Caught by} \\
\midrule
Grader mis-imported under a wrong root & ``self-training does not compound'' & identity-kernel assertion \\
Grading device index absent on a 1--2 GPU allocation & ``the model writes bad kernels'' & same assertion, run \emph{on the allocation shape the jobs use} \\
Raw generations graded without extracting the submission & ``the 14B teacher solves 0/29'' & non-zero coverage precondition \\
Score saturated at the correctness floor & ``lineage $\approx$ reset'' & score-histogram check \\
Compiled-baseline cache hitting a filesystem inode quota & ``4 of 6 rounds produced nothing'' & output-path/quota correlation \\
\bottomrule
\end{tabular}
\caption{Five defects, each of which independently manufactured a null. The middle column is
what the artifact would have claimed had the defect not been found.}
\label{tab:defects}
\end{table}

\paragraph{A clean zero is a bug hypothesis before it is a finding.}
The operational signature that caught most of these is distributional, not statistical. A
\emph{real} null is noisy: it has per-task spread, partial successes, seed-to-seed variance. A
null in which every task fails identically, or every score takes the same value, is
instrumentation. We therefore treat suspicious tidiness as evidence \emph{against} the result,
inverting the usual intuition that clean data is good data.

\paragraph{Correctness of the instrument is not sufficient; it must also have range.}
This is the failure mode we consider most transferable, because it defeats the obvious check.
Our headroom-normalised score is
$0.5 + 0.5\cdot\mathrm{clamp}\!\left(\frac{sp-1}{\mathrm{ceiling}-1}\right)$,
which is well-formed, was verified end-to-end on known-good inputs, and was working exactly as
specified. It was nonetheless dead as a measurement: on H100, \textbf{76\% of all scores across
every run and round were exactly $0.50$} --- precisely the value of a correct kernel at
parity with the baseline. Models reach the correctness floor in one to three rounds and stop,
because essentially no model from 0.5B to 14B produces a kernel faster than \texttt{torch.compile}
on this hardware.\footnote{The 14B teacher's best speedups across the bank: nine tasks at
$1.00\times$, two at $1.08\times$, and one each at $1.05\times$, $1.01\times$, $2.03\times$.}
With both arms pinned to the same absorbing value, $\textsf{lineage}-\textsf{reset}$ is forced
to zero \emph{by construction}, and the experiment could not have detected compounding in
either direction.

Critically, this is not repairable by task curation. Filtering the bank against an
H100 anchor with a $1.3\times$ minimum roofline headroom kept 23 of 29 tasks and still returned
a frozen-base mean best score of $0.511$, with every surviving task at $\approx 0.50$ and every
L3 task at $\textrm{correct\_rate}=0$. The difficulty distribution is bimodal with no middle
band: tasks are either unreachable or solved-at-parity. A filter selects among existing tasks;
it cannot create a regime in which the model is correct \emph{and} has headroom it can capture.

\paragraph{The estimator, not the normalisation, is the thing to fix.}
A saturated metric invites the wrong repair. The instinct is to re-normalise; the actual
problem is that best-of-$k$ over a speed score is the wrong \emph{estimator} for a population
sitting at the correctness floor. Best-of-$k$ answers ``can it ever,'' and returns an identical
$0.500$ whether a model solves a task once in eight attempts or eight times out of eight ---
two plainly different models. Where the live axis is reliability rather than speed, pass-rate
over $k$ separates those cases by a factor of eight in resolution.

We report such runs under a separate metric (pre-registration, Amendment~1), never pooled with the
headroom leaderboards, and with two properties stated up front: the denominator is $k$, so a
model emitting nothing parseable scores zero rather than being excluded; and it does not reward
speed, so it tests a \emph{weaker and different} claim than the headroom score. It was also
adopted \emph{after} observing saturation --- a forking path that we record as such rather than
present as the original design.

\paragraph{Check the denominator of the mechanism, not just the metric.}
The same failure shape recurs wherever a rate is computed over events that never occurred. Two
frontier models in our procedure-RSI track recorded $n_{\textrm{strategies}}=0$ in every round:
their output was never parsed into a strategy list, so the self-modification channel never
engaged. The resulting $-0.234$ ``self-degradation'' was widely quoted from an earlier draft of
this work and is a parse failure. Likewise, our self-play primary metric requires the live
author to produce accepted valid tasks; demanding at least five qualifies only 6 of 24 open and
5 of 11 closed runs. \textbf{A metric computed over zero events is undefined, not zero}, and the
distinction is invisible in the aggregate.

\paragraph{What we propose as standard practice.}
We ship these as executable gates rather than advice, and we recommend them for any
self-improvement benchmark reporting a null:
\begin{enumerate}[leftmargin=1.6em,itemsep=1pt]
\item \textbf{Known-good input, known-good output}, asserted before every batch and run
\emph{on the same allocation shape as the experiment} --- one of our defects lived entirely in
the difference between a 1-GPU and a 3-GPU allocation.
\item \textbf{A dynamic-range precondition}: demonstrate the primary metric takes at least two
distinguishable values on this hardware and this task bank, \emph{before} interpreting a
contrast computed from it.
\item \textbf{A positive control on the causal path}: an intervention known to move the primary
metric (for us, injecting verified solutions on tasks the student failed). A null is
interpretable only once the harness has been shown to register a positive.
\item \textbf{Mechanism denominators reported alongside every rate}, so undefined is
distinguishable from zero.
\item \textbf{Automated cross-artifact consistency checking}: every number in the manuscript
re-derived from the data files at build time. Ours catches stale figures and claim strings that
contradict the data they cite --- both of which occurred.
\end{enumerate}

We emphasise that this is not a catalogue of unusual carelessness. Each defect was a single
line, each was locally reasonable, and each was caught only because an assertion was looking
for it. The generalisable claim is structural: \textbf{in a benchmark whose headline is a
difference between two arms of one pipeline, the cheapest result to produce accidentally is a
null}, and the field currently has no convention requiring authors to rule that out.
